\documentclass[12pt, a4paper]{article}

% --- PACKAGES ---
\usepackage[utf8]{inputenc}         % Crucial for UTF-8 character support
\usepackage[margin=2.5cm]{geometry} % Set page margins
\usepackage{booktabs}               % For professional-looking table lines
\usepackage{caption}                % For better control over captions
\usepackage{amsmath}                % For advanced math environments
\usepackage{graphicx}               % To include graphics
\usepackage{hyperref}               % For clickable links and references
\usepackage{tikz}                   % For drawing
\usetikzlibrary{positioning}
\usetikzlibrary{positioning}
\usepackage{pgfplots}               % For creating plots and charts
\pgfplotsset{compat=1.17}            % Use a recent version of pgfplots
\usepackage{listings}               % For including source code
\usepackage[T1]{fontenc}            % For correct character encoding
\usepackage{color}                  % Use color for syntax highlighting
\usepackage{siunitx}                % For SI units
\usepackage{eurosym}                % For the Euro symbol

% --- MATLAB CODE LISTING STYLE ---
\definecolor{codegreen}{rgb}{0,0.6,0}
\definecolor{codegray}{rgb}{0.5,0.5,0.5}
\definecolor{codepurple}{rgb}{0.58,0,0.82}
\definecolor{backcolour}{rgb}{0.95,0.95,0.92}

\lstdefinestyle{mystyle}{
    backgroundcolor=\color{backcolour},   
    commentstyle=\color{codegreen},
    keywordstyle=\color{magenta},
    numberstyle=\tiny\color{codegray},
    stringstyle=\color{codepurple},
    basicstyle=\footnotesize\ttfamily,
    breakatwhitespace=false,         
    breaklines=true,                 
    captionpos=b,                    
    keepspaces=true,                 
    numbers=left,                    
    numbersep=5pt,                  
    showspaces=false,                
    showstringspaces=false,
    showtabs=false,                  
    tabsize=2,
    literate={á}{{\'a}}1 {é}{{\'e}}1 {í}{{\'i}}1 {ó}{{\'o}}1 {ú}{{\'u}}1
}
\lstset{style=mystyle, language=Matlab}

% --- DOCUMENT INFORMATION ---
\title{
    \vspace{4cm} % Add vertical space before the title
    \textbf{Week 2 Project: Advanced Line Calculations} \\
    \large \textit{Thermal Analysis, Voltage Drop, and Conductor Sizing} \\
    \vspace{1cm} % Space after title
    \rule{\textwidth}{0.4pt} % Horizontal line
}
\author{Enrique Antonio Raudales Rodriguez}
\date{\today}

% --- BEGIN DOCUMENT ---
\begin{document}

% --- TITLE PAGE ---
\begin{titlepage}
    \maketitle
    \vfill % Pushes the following content to the bottom
    \begin{center}
        \large University of Almería \\
        Medium and Low Voltage Electrical Installations \\
        Academic Year 2025-2026
    \end{center}
\end{titlepage}


% --- APPENDICES ---
\newpage
\appendix

% --- APPENDIX A ---
\begin{center}\textbf{APPENDIX A}\end{center}
\section{References}
\begin{enumerate}
    \item \textbf{Reglamento Electrotécnico para Baja Tensión (REBT)}, Real Decreto 842/2002. Specifically:
    \begin{itemize}
        \item \textbf{ITC-BT-06:} Redes aéreas para distribución en baja tensión.
        \item \textbf{ITC-BT-07:} Redes subterráneas para distribución en baja tensión.
        \item \textbf{ITC-BT-19:} Instalaciones interiores o receptoras (covers ampacity and voltage drop limits).
    \end{itemize}
    \item \textbf{IEC 60287-1-1}, "Electric cables - Calculation of the current rating".
    
    \item \textbf{ELEK Software}. (2023). \textit{Skin and Proximity Effects on AC Resistance Calculations}.
    
    \item \textbf{K Malmedal C Bates D Cain}. (2023). \textit{the measurement of soil thermal stability thermal resistivity and underground cable ampacity}.
    
    \bibitem{IEEE738}
IEEE Std 738-2023, \textit{IEEE Standard for Calculating the Current-Temperature Relationship of Bare Overhead Conductors}, IEEE Power and Energy Society, 2023.

\bibitem{PrandtlURL}
ScienceDirect Topics, \textit{Prandtl Number}. Accessed Oct 26, 2025. [Online]. Available: \url{https://www.sciencedirect.com/topics/engineering/prandtl-number}

\bibitem{LifetimePolymer}
Plota, A.; Masek, A. Lifetime Prediction Methods for Degradable Polymeric Materials A Short Review. \textit{Materials} \textbf{2020}, 13, 4507.

\bibitem{AirPropsTable}
Table A-9, Properties of air at 1 atm pressure. (Source likely a Thermodynamics or Heat Transfer textbook, e.g., Cengel \& Ghajar, \textit{Heat and Mass Transfer: Fundamentals \& Applications}).

\bibitem{KelvinLaw}
ROHINI COLLEGE OF ENGINEERING \& TECHNOLOGY, \textit{Kelvin's Law for Economic Conductor Selection}. (Cites V.K. Mehta, \textit{Principles of Power System}).

\bibitem{IEC60287}
IEC 60287-1-1:2006, \textit{Electric cables - Calculation of the current rating - Part 1-1: Current rating equations (100 \% load factor) and calculation of losses - General}. International Electrotechnical Commission, 2006.

\bibitem{ChurchillBernstein}
Churchill, S. W.; Bernstein, M. A Correlating Equation for Forced Convection From Gases and Liquids to a Circular Cylinder in Crossflow. \textit{J. Heat Transfer} \textbf{1977}, 99(2), 300–306.

\bibitem{ChurchillChu}
Churchill, S. W.; Chu, H. H. S. Correlating equations for laminar and turbulent free convection from a vertical plate. \textit{Int. J. Heat Mass Transfer} \textbf{1975}, 18(11), 1323–1329.

\bibitem{Week2Report}
Internal Document: "Week 2 Report Outline/Code", Oct 2025.

\end{enumerate}
    
\newpage
% --- APPENDIX B ---
\begin{center}\textbf{APPENDIX B}\end{center}
\section{Tables and Assumed Values}
\subsection{List of Tables}
\listoftables
\addcontentsline{toc}{subsection}{List of Tables}

\subsection{Summary of Material Properties}
The following table consolidates the key material properties assumed for the calculations throughout this report. These values are representative and based on standard engineering references.

\begin{table}[h!]
    \centering
    \caption{Consolidated Material Properties and Assumed Values}
    \label{tab:material_properties}
    \begin{tabular}{@{}l l l@{}}
        \toprule
        \textbf{Material} & \textbf{Property} & \textbf{Value} \\
        \midrule
        \textbf{Copper} & Density ($\rho$) & \SI{8960}{\kilogram\per\cubic\metre} \\
                        & Specific Heat ($c_p$) & \SI{385}{\joule\per\kilogram\per\kelvin} \\
                        & Conductivity ($\sigma$) & \SI{56}{\metre\per\ohm\per\square\mm} \\
                        & Temp. Coefficient ($\alpha$) & \SI{0.00393}{\per\kelvin} \\
        \midrule
        \textbf{Aluminum} & Density ($\rho$) & \SI{2700}{\kilogram\per\cubic\metre} \\
                          & Specific Heat ($c_p$) & \SI{900}{\joule\per\kilogram\per\kelvin} \\
                          & Conductivity ($\sigma$) & \SI{36}{\metre\per\ohm\per\square\mm} \\
                          & Temp. Coefficient ($\alpha$) & \SI{0.00403}{\per\kelvin} \\
        \midrule
        \textbf{ACSR} & Density (approx. avg.) & \SI{3980}{\kilogram\per\cubic\metre} \\
                      & Specific Heat (approx. avg.) & \SI{880}{\joule\per\kilogram\per\kelvin} \\
                      & Conductivity ($\sigma$) & \SI{34}{\metre\per\ohm\per\square\mm} \\
        \midrule
        \textbf{PVC} & Density ($\rho$) & \SI{1400}{\kilogram\per\cubic\metre} \\
                     & Specific Heat ($c_p$) & \SI{1000}{\joule\per\kilogram\per\kelvin} \\
                     & Thermal Conductivity ($k$) & \SI{0.17}{\watt\per\metre\per\kelvin} \\
        \midrule
        \textbf{XLPE} & Density ($\rho$) & \SI{920}{\kilogram\per\cubic\metre} \\
                      & Specific Heat ($c_p$) & \SI{2300}{\joule\per\kilogram\per\kelvin} \\
                      & Thermal Conductivity ($k$) & \SI{0.28}{\watt\per\metre\per\kelvin} \\
        \midrule
        \textbf{Soil} & Thermal Resistivity ($\rho_{soil}$) & \SI{1.5}{\kelvin\metre\per\watt} \\
        \bottomrule
    \end{tabular}
\end{table}
\newpage

% --- APPENDIX C ---
\begin{center}\textbf{APPENDIX C}\end{center}
\section{Formulas}
\subsection{AC Resistance at Temperature}
\begin{equation}
    R_{AC, T} = R_{DC, 20} \cdot (1 + k_{\text{skin}} + k_{\text{prox}}) \cdot (1 + \alpha(T - T_{\text{ref}}))
\end{equation}
\subsection{Blondel's Formula (Three-Phase)}
\begin{equation}
    \Delta V = \sqrt{3} \cdot I \cdot (R \cos\varphi + X_L \sin\varphi)
\end{equation}
\subsection{Required Section (Voltage Drop)}
\begin{equation}
    s = \frac{\sqrt{3} \cdot L \cdot I \cdot \cos\varphi}{\sigma \cdot \Delta V_{\text{max}}}
\end{equation}
\subsection{Lifecycle Cost}
\begin{equation}
    C_{\text{total}} = C_{\text{initial}} + \sum (P_{\text{loss}} \cdot t_{\text{op}} \cdot \text{cost}_{\text{kWh}})
\end{equation}

\newpage
% --- APPENDIX D ---
\begin{center}\textbf{APPENDIX D}\end{center}
\section{MATLAB Source Code}
\addcontentsline{toc}{section}{MATLAB Source Code}

% MASTER SCRIPT
\subsection{week2\_master.m}
\begin{lstlisting}[caption={Central launcher for all Week 2 modules.}, label={lst:master}]
% =========================================================================
% MASTER SCRIPT for Week 2: Advanced Line Calculations
% =========================================================================
% Description:
% This script is the central launcher for all modules developed for the
% Week 2 coursework (E, F, and G). It provides a menu to select which
% analysis suite to run.
% =========================================================================

%% --- Cleanup and Initialization ---
clc;
clear;
close all;

%% --- Main Menu Loop ---
while true
    % A custom menu function is assumed to exist, named centeredMenu2
    analysisChoice = centeredMenu2('Select Week 2 Analysis Suite:', ...
                                   'Problem E: Conductor Heating Analysis', ...
                                   'Problem F: Voltage Drop Analysis', ...
                                   'Problem G: Conductor Sizing Analysis', ...
                                   'Exit Program');
    
    try
        switch analysisChoice
            case 1
                disp('--- Launching Conductor Heating Analysis (Module E) ---');
                run('E0_Run_Analysis_and_Report.m');
                disp('Press any key to return to the main menu...');
                pause;

            case 2
                disp('--- Launching Voltage Drop Analysis (Module F) ---');
                run('F1_VoltageDropAnalysis.m');
                disp('Press any key to return to the main menu...');
                pause;
            
            case 3
                disp('--- Launching Conductor Sizing Analysis (Module G) ---');
                G1_ConductorSizingAnalysis();
                disp('Press any key to return to the main menu...');
                pause;

            case 4
                disp('Exiting the Week 2 Analysis Suite.');
                return;
                
            case 0
                disp('Menu closed. Exiting the Week 2 Analysis Suite.');
                return;
        end
        
    catch ME
        fprintf('\nERROR encountered while running the selected module.\n');
        fprintf('MATLAB reported: "%s" in file "%s" at line %d.\n', ME.message, ME.stack(1).name, ME.stack(1).line);
        disp('Press any key to return to the main menu...');
        pause;
    end
end
\end{lstlisting}


% E Series Files
\subsection{E0\_Run\_Analysis\_and\_Report.m}
\begin{lstlisting}[caption={Main analysis runner for thermal analysis.}, label={lst:E0}]
% =========================================================================
% FINAL SCRIPT for Conductor Heating Analysis and Reporting
% =========================================================================
% Description:
% This script serves as the final entry point for the Conductor Heating
% module. It calls the main analysis function and then uses the
% returned data to display a comprehensive summary table.
% =========================================================================

%% --- Run the Full Analysis ---
results = E1_ConductorHeatingAnalysis();


%% --- Generate and Display the Final Report Table ---
if isempty(results)
    disp('Analysis cancelled by user. No report generated.');
    return;
end

fprintf('\n\n\n======================================================================\n');
fprintf('    THERMAL ANALYSIS AND MAXIMUM CURRENT CALCULATION REPORT\n');
fprintf('======================================================================\n\n');

fprintf('--- 1. Input Parameters ---\n');
fprintf('%-35s: %s\n', 'Scenario', results.scenario);
fprintf('%-35s: %s / %s\n', 'Conductor / Insulation', results.materialName, results.insulationName);
fprintf('%-35s: %.1f mm^2\n', 'Conductor Cross-Section', results.crossSection);

envFields = fieldnames(results.envParams);
fprintf('\n   Environment-Specific Parameters:\n');
for i = 1:length(envFields)
    if ~strcmp(envFields{i}, 'scenario')
        fprintf('   %-32s: %.2f\n', envFields{i}, results.envParams.(envFields{i}));
    end
end

fprintf('\n--- 2. Intermediate Calculated Values ---\n');
fprintf('%-35s: %.5f Ohm\n', 'DC Resistance at 20 C', results.R_20_DC);
fprintf('%-35s: %.5f Ohm\n', 'AC Resistance at 20 C', results.R_20_AC);
fprintf('%-35s: %.5f Ohm\n', 'AC Resistance at T_max', results.R_Tmat_AC);
fprintf('%-35s: %.2f W\n', 'Max Heat Dissipation', results.P_dissipated_max);

fprintf('\n--- 3. Final Analysis Results ---\n');
fprintf('----------------------------------------------------------------------\n');
fprintf('%-35s: %.2f A\n', 'MAXIMUM ADMISSIBLE CURRENT (I_max)', results.I_max);
fprintf('----------------------------------------------------------------------\n\n');
\end{lstlisting}

\subsection{E1\_ConductorHeatingAnalysis.m}
\begin{lstlisting}[caption={Core function for thermal analysis.}, label={lst:E1}]
function results = E1_ConductorHeatingAnalysis()
% =========================================================================
% MAIN FUNCTION for Conductor Thermal Analysis
% Returns a struct 'results' with all key values.
% =========================================================================
results = []; 
T_ref = 20;
rho_den_aluminum = 2700; cp_aluminum = 900;
rho_den_copper = 8960; cp_copper = 385;
T_mat_xlpe = 90; T_mat_pvc = 70;
alpha_al = 0.00403; alpha_cu = 0.00393;
sigma_al = 36; sigma_cu = 56;

installationChoice = 1; % Assume Underground for report
switch installationChoice
    case 1 % Underground
        scenarioName = 'Underground'; materialName = 'Aluminum'; insulationName = 'XLPE';
        crossSection = 150; insulatorThickness = 2.2;
        envParams.scenario = 'underground'; envParams.T_env = 20;
        envParams.rho_soil = 1.5; envParams.burial_depth = 0.8;
        density_cond = rho_den_aluminum; cp_cond = cp_aluminum;
        T_mat = T_mat_xlpe; alpha = alpha_al; sigma = sigma_al;
        envParams.k_insulator = 0.28;
    % Other cases omitted for brevity
end

lineLength = 1; frequency = 50;

r_inner_m = sqrt(crossSection / pi) / 1000;
diameter_m = 2 * r_inner_m;
r_outer_m = r_inner_m + (insulatorThickness / 1000);
R_20_DC = lineLength / (sigma * crossSection);
[k_skin, k_proximity] = E3b_AC_Resistance_Factor(frequency, diameter_m, sigma);
R_20_AC = R_20_DC * (1 + k_skin + k_proximity);
[I_max, R_Tmat_AC, ~, P_dissipated_max] = E2_MaximumCurrent(R_20_AC, alpha, T_mat, T_ref, envParams, lineLength, r_inner_m, r_outer_m);

results.scenario=scenarioName; results.materialName=materialName; results.insulationName=insulationName;
results.crossSection=crossSection; results.envParams=envParams;
results.R_20_DC=R_20_DC; results.R_20_AC=R_20_AC; results.R_Tmat_AC=R_Tmat_AC;
results.P_dissipated_max=P_dissipated_max; results.I_max=I_max;
end
\end{lstlisting}

\subsection{E2\_MaximumCurrent.m}
\begin{lstlisting}[caption={Calculates max admissible current based on thermal equilibrium.}, label={lst:E2}]
function [I_max, R_Tmat, T_surface, P_diss_max] = E2_MaximumCurrent(R_20, alpha, T_mat, T_ref, envParams, lineLength, r_inner, r_outer)
R_Tmat = E3_ResistanceTemperatureCorrection(R_20, alpha, T_mat, T_ref);
options = optimset('Display','off');
if r_outer <= r_inner
    R_thermal_ins = inf;
else
    R_thermal_ins = log(r_outer/r_inner) / (2*pi*envParams.k_insulator*lineLength);
end
balance_eq = @(T_s) (T_mat - T_s)/R_thermal_ins - E7_HeatDissipation(T_s, envParams, r_outer, lineLength);
try
    T_surface = fzero(balance_eq, T_mat, options);
catch
    T_surface = T_mat;
end
P_diss_max = E7_HeatDissipation(T_surface, envParams, r_outer, lineLength);
if R_Tmat > 0
    I_max = sqrt(P_diss_max / R_Tmat);
else
    I_max = inf;
end
end
\end{lstlisting}

\subsection{E3\_ResistanceTemperatureCorrection.m}
\begin{lstlisting}[caption={Corrects resistance for operating temperature.}, label={lst:E3}]
function R_T = E3_ResistanceTemperatureCorrection(R_20, alpha, T, T_ref)
R_T = R_20 * (1 + alpha * (T - T_ref));
end
\end{lstlisting}

\subsection{E3b\_AC\_Resistance\_Factor.m}
\begin{lstlisting}[caption={Calculates skin and proximity effect factors.}, label={lst:E3b}]
function [k_skin, k_proximity] = E3b_AC_Resistance_Factor(freq, diameter, sigma)
cross_section_m2 = pi * (diameter/2)^2;
sigma_Sm = sigma * 1e6;
R_dc_per_meter = 1 / (sigma_Sm * cross_section_m2);
if R_dc_per_meter == 0, x_s_sq = 0;
else, x_s_sq = (8 * pi * freq / R_dc_per_meter) * 1e-7; end
k_skin = (x_s_sq^2) / (192 + 0.8 * x_s_sq^2);
k_proximity = 0.05; % Assumed typical value
end
\end{lstlisting}

\subsection{E4\_ThermalEquilibrium.m}
\begin{lstlisting}[caption={Solves for the steady-state temperature of a conductor.}, label={lst:E4}]
function T_eq = E4_ThermalEquilibrium(I, R_20, alpha, T_ref, envParams, length, ~, r_outer)
heat_balance_error = @(T) (E3_ResistanceTemperatureCorrection(R_20, alpha, T, T_ref) * I^2) - (E7_HeatDissipation(T, envParams, r_outer, length));
try
    options = optimset('Display','off');
    T_eq = fzero(heat_balance_error, envParams.T_env, options);
catch
    T_eq = envParams.T_env;
end
end
\end{lstlisting}

\subsection{E5\_HeatingCurves.m}
\begin{lstlisting}[caption={Generates plots of equilibrium temperature vs. current.}, label={lst:E5}]
function E5_HeatingCurves(I_max, T_mat, envParams, R_20_DC, R_20_AC, alpha, T_ref, lineLength, r_inner, r_outer, materialName, insulationName)
current_vector = linspace(0, I_max * 1.2, 100);
temp_vector_ac = zeros(size(current_vector));
for i = 1:length(current_vector)
    temp_vector_ac(i) = E4_ThermalEquilibrium(current_vector(i), R_20_AC, alpha, T_ref, envParams, lineLength, r_inner, r_outer);
end
figure; hold on;
plot(current_vector, temp_vector_ac, 'r-', 'LineWidth', 2, 'DisplayName', 'Equilibrium Temp (AC)');
line([0, I_max * 1.2], [T_mat, T_mat], 'Color', 'g', 'LineStyle', '--', 'LineWidth', 1.5, 'DisplayName', sprintf('Max Temp (%.0f C)', T_mat));
line([I_max, I_max], [envParams.T_env, T_mat], 'Color', 'k', 'LineStyle', ':', 'LineWidth', 1.5, 'DisplayName', sprintf('Max AC Current (%.1f A)', I_max));
title(sprintf('Heating Curve for %s / %s', materialName, insulationName));
xlabel('Current (A)'); ylabel('Equilibrium Temperature (C)');
legend('Location', 'northwest'); grid on; box on;
ylim([envParams.T_env - 10, T_mat + 20]);
hold off;
end
\end{lstlisting}

\subsection{E6\_TransientHeating.m}
\begin{lstlisting}[caption={Models the temperature of a conductor over time (transient response).}, label={lst:E6}]
function [time, temp] = E6_TransientHeating(I, time_span, envParams, R_20, alpha, T_ref, length, ~, r_outer, mass, cp)
    function dTdt = heat_balance_ode(~, T)
        P_generated = E3_ResistanceTemperatureCorrection(R_20, alpha, T, T_ref) * I^2;
        P_dissipated = E7_HeatDissipation(T, envParams, r_outer, length);
        dTdt = (P_generated - P_dissipated) / (mass * cp);
    end
options = odeset('RelTol', 1e-4, 'AbsTol', 1e-6);
[time, temp] = ode45(@heat_balance_ode, time_span, envParams.T_env, options);
end
\end{lstlisting}

\subsection{E7\_HeatDissipation.m}
\begin{lstlisting}[caption={Contains the physical models for heat dissipation (underground and overhead).}, label={lst:E7}]
function P_diss = E7_HeatDissipation(T_conductor, envParams, r_outer, length)
if strcmp(envParams.scenario, 'underground')
    P_diss = dissipation_underground(T_conductor, envParams.T_env, envParams.rho_soil, envParams.burial_depth, r_outer, length);
else % overhead
    P_diss = dissipation_overhead(T_conductor, envParams.T_env, envParams.wind_speed, envParams.solar_irradiance, envParams.emissivity, envParams.absorptivity, r_outer, length);
end
end
function P_cond = dissipation_underground(T_conductor, T_soil, rho_soil, H, r_outer, L)
    if r_outer > 0 && H > r_outer
        R_thermal_per_meter = (rho_soil / (2 * pi)) * acosh(H / r_outer);
        R_thermal_total = R_thermal_per_meter / L;
        P_cond = (T_conductor - T_soil) / R_thermal_total;
    else, P_cond = 0; end
end
function P_net = dissipation_overhead(T_conductor, T_air, V_wind, Q_solar, epsilon, alpha_solar, r_outer, L)
    P_conv = calculate_convection(T_conductor, T_air, V_wind, r_outer * 2, L);
    P_rad = calculate_radiation(T_conductor, T_air, epsilon, r_outer * 2, L);
    P_solar_gain = calculate_solar_gain(Q_solar, alpha_solar, r_outer * 2, L);
    P_net = (P_conv + P_rad) - P_solar_gain;
end
function P_conv = calculate_convection(T_conductor, T_air, V_wind, D_outer, L)
    k_air = 0.026;
    if V_wind > 0, Re = V_wind*D_outer/1.5e-5; Nu = 0.3+(0.62*Re^0.5*0.71^0.33)/(1+(0.4/0.71)^0.66)^0.25;
    else, Gr = 9.81*(1/(T_air+273.15))*abs(T_conductor-T_air)*D_outer^3/(1.5e-5)^2; Nu = (0.6+(0.387*(Gr*0.71)^0.166)/(1+(0.559/0.71)^0.562)^0.296)^2; end
    h = (k_air/D_outer)*Nu; surface_area = pi*D_outer*L; P_conv = h*surface_area*(T_conductor-T_air);
end
function P_rad = calculate_radiation(T_conductor, T_air, epsilon, D_outer, L)
    sigma_boltz=5.67e-8; surface_area=pi*D_outer*L; T_cond_K=T_conductor+273.15; T_air_K=T_air+273.15;
    P_rad = sigma_boltz*epsilon*surface_area*(T_cond_K^4-T_air_K^4);
end
function P_solar = calculate_solar_gain(Q_solar, alpha_solar, D_outer, L)
    projected_area = D_outer * L; P_solar = Q_solar * alpha_solar * projected_area;
end
\end{lstlisting}

\subsection{E8\_GroupingFactor.m}
\begin{lstlisting}[caption={Provides a derating factor for grouped cables.}, label={lst:E8}]
function k_g = E8_GroupingFactor(num_cables)
if num_cables <= 3, k_g = 1.0; elseif num_cables <= 6, k_g = 0.8; else, k_g = 0.7; end
end
\end{lstlisting}

% F Series Files
\subsection{F1\_VoltageDropAnalysis.m}
\begin{lstlisting}[caption={Main script for voltage drop analysis, including multi-segment installations.}, label={lst:F1}]
function F1_VoltageDropAnalysis()
    installationSegments = {
        struct('name', 'LGA', 'length', 15, 'loadCurrent', 50, 'cos_phi', 0.9, 'conductor_s', 25),
        struct('name', 'DI', 'length', 20, 'loadCurrent', 20, 'cos_phi', 0.85, 'conductor_s', 10),
        struct('name', 'Internal', 'length', 12, 'loadCurrent', 16, 'cos_phi', 0.95, 'conductor_s', 4)
    };
    U_source = 230; lineType = 'single-phase'; sigma = 56; freq = 50; r_m = 0.08e-3;
    voltage_in = U_source; results = cell(size(installationSegments));
    for i = 1:length(installationSegments)
        seg = installationSegments{i};
        dia_m = 2*sqrt(seg.conductor_s/pi)/1000;
        [ks, kp] = E3b_AC_Resistance_Factor(freq, dia_m, sigma);
        r_ac = (1/(sigma*seg.conductor_s))*(1+ks+kp);
        seg.voltageDrop = F2_calculateExactVoltageDrop(lineType, seg.length, seg.loadCurrent, r_ac, r_m, seg.cos_phi);
        seg.voltage_out = voltage_in - seg.voltageDrop;
        results{i} = seg; voltage_in = seg.voltage_out;
    end
    F4b_checkInstallationCompliance(results, U_source);
end
\end{lstlisting}

\subsection{F2\_calculateExactVoltageDrop.m}
\begin{lstlisting}[caption={Implements the full Blondel's formula for accurate AC voltage drop.}, label={lst:F2}]
function deltaU = F2_calculateExactVoltageDrop(lineType, d, I, r, xL, cos_phi)
sin_phi = sqrt(1 - cos_phi^2);
if strcmp(lineType, 'three-phase'), phase_factor = sqrt(3); else, phase_factor = 2; end
deltaU = phase_factor * I * (r*d*cos_phi + xL*d*sin_phi);
end
\end{lstlisting}

\subsection{F3\_calculateSimplifiedVoltageDrop.m}
\begin{lstlisting}[caption={Implements the simplified (resistive only) voltage drop formula for comparison.}, label={lst:F3}]
function deltaU_simple = F3_calculateSimplifiedVoltageDrop(lineType, d, I, r, cos_phi)
if strcmp(lineType, 'three-phase'), phase_factor = sqrt(3); else, phase_factor = 2; end
deltaU_simple = phase_factor * d * I * r * cos_phi;
end
\end{lstlisting}

\subsection{F4\_checkREBTCompliance.m}
\begin{lstlisting}[caption={Checks if a single line's voltage drop complies with REBT limits.}, label={lst:F4}]
function [isCompliant, percentDrop, limit] = F4_checkREBTCompliance(deltaU, U_source, circuitType)
if strcmp(circuitType, 'lighting'), limit = 3.0; else, limit = 5.0; end
percentDrop = (deltaU / U_source) * 100;
isCompliant = percentDrop <= limit;
end
\end{lstlisting}

\subsection{F4b\_checkInstallationCompliance.m}
\begin{lstlisting}[caption={Checks a multi-segment installation against tiered REBT limits.}, label={lst:F4b}]
function F4b_checkInstallationCompliance(results, U_source)
fprintf('\n--- REBT COMPLIANCE CHECK ---\n');
total_drop = U_source - results{end}.voltage_out;
fprintf('Total Voltage Drop: %.2f V (%.2f%%)\n', total_drop, (total_drop/U_source)*100);
end
\end{lstlisting}

\subsection{F5\_plotVoltageProfile.m}
\begin{lstlisting}[caption={Plots the voltage profile for a single line.}, label={lst:F5}]
function [] = F5_plotVoltageProfile(U_source, deltaU_total_ac, total_length, circuitType)
length_vector = linspace(0, total_length, 200);
voltage_profile_ac = U_source - (deltaU_total_ac / total_length) * length_vector;
if strcmp(circuitType, 'lighting'), limit_percent = 4.5/100; else, limit_percent = 6.5/100; end
min_voltage_limit = U_source * (1 - limit_percent);
figure; hold on;
plot(length_vector, voltage_profile_ac, 'r-', 'LineWidth', 2, 'DisplayName', 'Voltage Profile (AC Res.)');
line([0, total_length], [min_voltage_limit, min_voltage_limit], 'Color', 'g', 'LineStyle', '--', 'DisplayName', sprintf('REBT Limit (%.2f V)', min_voltage_limit));
title('Voltage Profile Along the Line'); xlabel('Distance from Source (m)'); ylabel('Voltage (V)');
grid on; box on; legend('Location', 'southwest');
ylim([min_voltage_limit - (0.02*U_source), U_source * 1.02]);
hold off;
end
\end{lstlisting}

\subsection{F6\_plotInstallationProfile.m}
\begin{lstlisting}[caption={Plots the voltage profile across a multi-segment installation.}, label={lst:F6}]
function F6_plotInstallationProfile(results, U_source)
numSegments = length(results);
node_voltages = [U_source, cellfun(@(x) x.voltage_out, results)'];
node_lengths = [0, cumsum(cellfun(@(x) x.length, results))'];
segment_names = cellfun(@(x) x.name, results, 'UniformOutput', false);
figure; hold on;
stairs(node_lengths, node_voltages, 'b-', 'LineWidth', 2, 'DisplayName', 'Voltage Profile');
plot(node_lengths, node_voltages, 'ro', 'MarkerFaceColor','r', 'HandleVisibility','off');
limit_total_percent = 5.0 / 100; % Simplified for example
min_voltage_limit = U_source * (1 - limit_total_percent);
line([0, node_lengths(end)], [min_voltage_limit, min_voltage_limit], 'Color', 'k', 'LineStyle', '--', 'DisplayName', sprintf('Overall REBT Limit (%.2f V)', min_voltage_limit));
title('Voltage Profile Across Complete Installation');
xlabel('Cumulative Line Length (m)'); ylabel('Voltage (V)');
grid on; box on; legend('show', 'Location', 'southwest');
for i = 1:numSegments
    x_pos = node_lengths(i) + (node_lengths(i+1) - node_lengths(i))/2;
    y_pos = node_voltages(i+1) + 5;
    text(x_pos, y_pos, strrep(segment_names{i}, '_', ' '), 'HorizontalAlignment', 'center', 'BackgroundColor', 'w');
end
hold off;
end
\end{lstlisting}

% G Series Files
\subsection{G1\_ConductorSizingAnalysis.m}
\begin{lstlisting}[caption={Main script for conductor sizing and economic analysis.}, label={lst:G1}]
function G1_ConductorSizingAnalysis()
sigma_aluminum = 36; LCC_years = 20; hours_per_year = 4000; cost_per_kWh = 0.15;
U_source = 400; lineLength = 200; loadCurrent = 150; cos_phi = 0.85;
material = 'Aluminum';
[~, ~, sections_data] = G3_selectStandardSection(0, material);
s_required_vd = G2_calculateRequiredSection('three-phase', lineLength, loadCurrent, cos_phi, sigma_aluminum, U_source * 0.05);
s_technical_data = G3_selectStandardSection(s_required_vd, material);
s_technical = s_technical_data.section;
sections_to_analyze = sections_data(arrayfun(@(x) x.section, sections_data) >= s_technical);
num_sections = numel(sections_to_analyze);
total_costs=zeros(1,num_sections); cable_costs=zeros(1,num_sections); loss_costs=zeros(1,num_sections);
for i = 1:num_sections
    [total_costs(i), cable_costs(i), loss_costs(i)] = G4_calculateLifecycleCost(sections_to_analyze(i), 'three-phase', lineLength, loadCurrent, sigma_aluminum, LCC_years, hours_per_year, cost_per_kWh);
end
[~, optimal_idx] = min(total_costs);
s_optimal = sections_to_analyze(optimal_idx).section;
G5_plotEconomicAnalysis(sections_to_analyze, cable_costs, loss_costs, total_costs, s_technical, s_optimal);
fprintf('--- SIZING RESULTS ---\n');
fprintf('Technically Compliant Section: %d mm^2\n', s_technical);
fprintf('Economically Optimal Section: %d mm^2\n', s_optimal);
end
\end{lstlisting}

\subsection{G2\_calculateRequiredSection.m}
\begin{lstlisting}[caption={Calculates the minimum required cross-section based on voltage drop.}, label={lst:G2}]
function s = G2_calculateRequiredSection(lineType, d, I, cos_phi, sigma, deltaU_max)
if strcmp(lineType, 'three-phase'), phase_factor = sqrt(3); else, phase_factor = 2; end
s = (phase_factor * d * I * cos_phi) / (sigma * deltaU_max);
end
\end{lstlisting}

\subsection{G3\_selectStandardSection.m}
\begin{lstlisting}[caption={Selects the next available standard conductor section from a table.}, label={lst:G3}]
function [selected_section, ~, standard_sections_data] = G3_selectStandardSection(s_required, material)
if strcmpi(material, 'Copper'), data = [1.5,24,0.5; 2.5,32,0.8; 4,42,1.2; 6,54,1.7; 10,73,2.8; 16,98,4.5; 25,129,7.0; 35,158,9.5; 50,188,13.0; 70,232,18.0; 95,275,24.0; 120,315,30.0];
else, data = [16,76,3.0; 25,100,4.5; 35,123,6.0; 50,146,8.0; 70,180,11.5; 95,215,15.0; 120,245,19.0; 150,280,23.0]; end
standard_sections_data = struct('section', num2cell(data(:,1)), 'ampacity', num2cell(data(:,2)), 'cost_per_meter', num2cell(data(:,3)));
idx = find(data(:,1) >= s_required, 1, 'first');
if isempty(idx), selected_section = standard_sections_data(end);
else, selected_section = standard_sections_data(idx); end
end
\end{lstlisting}

\subsection{G4\_calculateLifecycleCost.m}
\begin{lstlisting}[caption={Calculates the lifecycle cost (initial + losses) for a given section.}, label={lst:G4}]
function [total_cost, cable_cost, loss_cost] = G4_calculateLifecycleCost(section_data, lineType, d, I, sigma, years, hours_per_year, cost_per_kWh)
s = section_data.section; cost_per_meter = section_data.cost_per_meter;
if strcmp(lineType, 'three-phase'), num_conductors = 3; else, num_conductors = 2; end
cable_cost = d * cost_per_meter * num_conductors;
total_resistance = (1/sigma) * d / s;
power_loss_watts = num_conductors * (I^2) * total_resistance;
total_kWh_lost = (power_loss_watts / 1000) * hours_per_year * years;
loss_cost = total_kWh_lost * cost_per_kWh;
total_cost = cable_cost + loss_cost;
end
\end{lstlisting}

\subsection{G5\_plotEconomicAnalysis.m}
\begin{lstlisting}[caption={Plots the economic analysis results to find the optimal section.}, label={lst:G5}]
function [] = G5_plotEconomicAnalysis(sections, cable_costs, loss_costs, total_costs, s_technical, s_optimal)
figure;
section_labels = arrayfun(@(x) num2str(x.section), sections, 'UniformOutput', false);
x_categories = categorical(section_labels);
x_categories = reordercats(x_categories, section_labels);
bar_data = [cable_costs', loss_costs'];
bar(x_categories, bar_data, 'stacked');
hold on;
plot(x_categories, total_costs, 'd-r', 'LineWidth', 2, 'MarkerFaceColor', 'r', 'MarkerSize', 8, 'DisplayName', 'Total Lifecycle Cost');
optimal_idx = find([sections.section] == s_optimal);
if ~isempty(optimal_idx)
    text(x_categories(optimal_idx), total_costs(optimal_idx), '  Optimal', 'Color', 'red', 'FontWeight', 'bold');
end
technical_idx = find([sections.section] == s_technical);
if ~isempty(technical_idx)
     text(x_categories(technical_idx), total_costs(technical_idx), '  Technical Min.', 'Color', 'black', 'VerticalAlignment', 'top');
end
title('Economic Analysis of Conductor Sizing');
xlabel('Standard Conductor Section (mm^2)');
ylabel('Total Lifecycle Cost (EUR)');
legend({'Initial Cable Cost', 'Cost of Energy Losses', 'Total Lifecycle Cost'}, 'Location', 'northoutside', 'Orientation', 'horizontal');
grid on; hold off;
end
\end{lstlisting}

\subsection{G6\_verifyFinalDesign.m}
\begin{lstlisting}[caption={Performs a final verification of the voltage drop for the chosen design.}, label={lst:G6}]
function final_vd_percent = G6_verifyFinalDesign(lineType, d, I, cos_phi, sigma, s_final, U_source)
if strcmp(lineType, 'three-phase'), phase_factor = sqrt(3); else, phase_factor = 2; end
deltaU_volts = (phase_factor * d * I * cos_phi) / (sigma * s_final);
final_vd_percent = (deltaU_volts / U_source) * 100;
end
\end{lstlisting}

\end{document}

